Power of P gives $PR\cdot PX=40^2-20^2=1200$, so $PX=48$. Use $O=(0,0),P=(20,0),Q=(40,0)$. Solving $OR=40$ and $PR=25$ gives $R=(\dfrac{275}{8},\dfrac{15\sqrt{119}}{8})$, up to reflection. Since X is on the opposite ray from P, its height is $-(\dfrac{48}{25})(\dfrac{15\sqrt{119}}{8})=-\dfrac{18\sqrt{119}}{5}$. With base $PQ=20$, the triangle area is $36\sqrt{119}$, hence $m+n=\boxed{155}$.
\begin{asy}
import graph;import geometry;usepackage("amsmath");size(7cm);
pair O=(0,0),P=(20,0),Q=(40,0),R=(275/8,15*sqrt(119)/8),X=P-48/25*(R-P);draw(circle(O,40));draw(O--Q);draw(R--X);draw(X--Q,blue);label("$O$",O,(-1,1));label("$P$",P,(0,1));label("$Q$",Q,(1,0));label("$R$",R,(1,1));label("$X$",X,(-1,-1));
\end{asy}

